\documentclass[../../main.tex]{subfiles}
\begin{document}
\section{向量场与光滑映射}\label{section:ch08-向量场与光滑映射}

若 $F\colon M\to N$ 是光滑映射，$X$ 是 $M$ 上的向量场，则对每一点 $p\in M$，通过将 $F$ 的微分作用于 $X_p$，得到向量 $\mathrm{d}F_p(X_p)\in T_{F(p)}N$。然而，这一般并不能在 $N$ 上定义一个向量场。例如，若 $F$ 不是满射，则无法决定给点 $q\in N\setminus F(M)$ 指派什么向量。若 $F$ 不是单射，则对 $N$ 的某些点，可能有通过在 $M$ 的不同点处对 $X$ 应用 $\mathrm{d}F$ 而得到的若干不同向量。

\begin{figure}[ht]
\centering
\includegraphics[width=0.8\textwidth]{fig-08-2.png}
\caption{微分可能不把向量场映到向量场}\label{figure:ch08-differential-vector-fields}
\end{figure}

假设 $F\colon M\to N$ 光滑，$X$ 是 $M$ 上的向量场，并假设恰好存在 $N$ 上的向量场 $Y$，使得对每个 $p\in M$ 有 $\mathrm{d}F_p(X_p)=Y_{F(p)}$。此时我们称向量场 $X$ 和 $Y$ 是\textbf{$F$-相关的}。下一个命题说明 $F$-相关向量场如何作用于光滑函数。

\begin{figure}[ht]
\centering
\includegraphics[width=0.8\textwidth]{fig-08-3.png}
\caption{$F$-相关向量场}\label{figure:ch08-f-related-vector-fields}
\end{figure}

\begin{proposition}\label{proposition:ch08-16}
设 $F\colon M\to N$ 是带边或不带边流形之间的光滑映射，$X\in\mathfrak{X}(M)$，$Y\in\mathfrak{X}(N)$。则 $X$ 和 $Y$ 是 $F$-相关的，当且仅当对定义在 $N$ 的开子集上的每个光滑实值函数 $f$，
\begin{align}
X(f\circ F)=(Yf)\circ F.\label{align:ch08-06}
\end{align}
\end{proposition}

\begin{proof}
对任意 $p\in M$ 和定义在 $F(p)$ 邻域内的任意光滑实值函数 $f$，
\begin{align*}
X(f\circ F)(p)=X_p(f\circ F)=\mathrm{d}F_p(X_p)f,
\end{align*}
而
\begin{align*}
(Yf)\circ F(p)=(Yf)\bigl(F(p)\bigr)=Y_{F(p)}f.
\end{align*}
因此，\eqref{align:ch08-06} 对所有 $f$ 成立，当且仅当对所有 $p$ 有 $\mathrm{d}F_p(X_p)=Y_{F(p)}$，即当且仅当 $X$ 和 $Y$ 是 $F$-相关的。
\end{proof}

\begin{example}\label{example:ch08-17}
设 $F\colon\mathbb{R}\to\mathbb{R}^2$ 为光滑映射 $F(t)=(\cos t,\sin t)$。则 $\mathrm{d}/\mathrm{d}t\in\mathfrak{X}(\mathbb{R})$ 与由
\begin{align*}
Y=x\frac{\partial}{\partial y}-y\frac{\partial}{\partial x}
\end{align*}
定义的向量场 $Y\in\mathfrak{X}(\mathbb{R}^2)$ 是 $F$-相关的。
\end{example}

\begin{exercise}\label{exercise:ch08-18}
用两种方式证明前一例子中的断言：直接由定义，以及使用\cref{proposition:ch08-16}。
\end{exercise}

重要的是要记住，对于给定的光滑映射 $F\colon M\to N$ 和向量场 $X\in\mathfrak{X}(M)$，$N$ 上可能不存在任何与 $X$ 是 $F$-相关的向量场。然而，有一种特殊情形总是存在这样的向量场，如下一个命题所示。

\begin{proposition}\label{proposition:ch08-19}
假设 $M$ 和 $N$ 是带边或不带边的光滑流形，$F\colon M\to N$ 是微分同胚。对每个 $X\in\mathfrak{X}(M)$，存在 $N$ 上唯一的光滑向量场与 $X$ 是 $F$-相关的。
\end{proposition}

\begin{proof}
$Y\in\mathfrak{X}(N)$ 与 $X$ 是 $F$-相关的，意味着对每个 $p\in M$ 有 $\mathrm{d}F_p(X_p)=Y_{F(p)}$。因此，若 $F$ 是微分同胚，则定义 $Y$ 为
\begin{align*}
Y_q=\mathrm{d}F_{F^{-1}(q)}\bigl(X_{F^{-1}(q)}\bigr).
\end{align*}
显然，如此定义的 $Y$ 是与 $X$ 是 $F$-相关的唯一（粗）向量场。注意 $Y\colon N\to TN$ 是如下光滑映射的复合：
\begin{align*}
N\xrightarrow{F^{-1}}M\xrightarrow{X}TM\xrightarrow{\mathrm{d}F}TN.
\end{align*}
因此 $Y$ 是光滑的。
\end{proof}

在前述命题的情形中，我们将与 $X$ 是 $F$-相关的唯一向量场记作 $F_*X$，并称为 $X$ 在 $F$ 下的\textbf{前推}。记住，只有当 $F$ 是微分同胚时 $F_*X$ 才有定义。\cref{proposition:ch08-19}的证明表明 $F_*X$ 由公式
\begin{align}
(F_*X)_q=\mathrm{d}F_{F^{-1}(q)}\bigl(X_{F^{-1}(q)}\bigr)\label{align:ch08-07}
\end{align}
显式定义。只要逆映射 $F^{-1}$ 可以显式计算，向量场的前推就可以直接由该公式计算。

\begin{example}\label{example:ch08-20}[计算向量场的前推]
设 $M$ 和 $N$ 是 $\mathbb{R}^2$ 的如下开子流形：
\begin{align*}
M=\{(x,y):y>0\text{ 且 }x+y>0\},\qquad
N=\{(u,v):u>0\text{ 且 }v>0\},
\end{align*}
并定义 $F\colon M\to N$ 为 $F(x,y)=(x+y,x/y+1)$。则 $F$ 是微分同胚，因为其逆容易计算：只需解 $(u,v)=(x+y,x/y+1)$ 求 $x$ 和 $y$，得到 $(x,y)=F^{-1}(u,v)=(u-u/v,u/v)$。我们来计算前推 $F_*X$，其中 $X$ 是 $M$ 上的如下光滑向量场：
\begin{align*}
X(x,y)=y^2\frac{\partial}{\partial x}\bigg|_{(x,y)}.
\end{align*}
$F$ 在点 $(x,y)\in M$ 处的微分由其 Jacobian 矩阵表示：
\begin{align*}
DF(x,y)=\begin{pmatrix}1&1\\ 1/y&-x/y^2\end{pmatrix},
\end{align*}
因此 $\mathrm{d}F_{F^{-1}(u,v)}$ 由矩阵
\begin{align*}
DF\left(u-\frac{u}{v},\frac{u}{v}\right)=\begin{pmatrix}1&1\\ v/u&-v^2/u\end{pmatrix}
\end{align*}
表示。对任意 $(u,v)\in N$，
\begin{align*}
X_{F^{-1}(u,v)}=\frac{u^2}{v^2}\frac{\partial}{\partial x}\bigg|_{F^{-1}(u,v)}.
\end{align*}
因此，在\eqref{align:ch08-07}中取 $p=(u,v)$，得到 $F_*X$ 的公式：
\begin{align*}
(F_*X)(u,v)=\frac{u^2}{v^2}\frac{\partial}{\partial u}\bigg|_{(u,v)}+\frac{u}{v}\frac{\partial}{\partial v}\bigg|_{(u,v)}.
\end{align*}
\end{example}

下一个推论直接由\cref{proposition:ch08-16}得到。

\begin{corollary}\label{corollary:ch08-21}
假设 $F\colon M\to N$ 是微分同胚，$X\in\mathfrak{X}(M)$。对任意 $f\in C^\infty(N)$，
\begin{align*}
(F_*X)f\circ F=X(f\circ F).
\end{align*}
\end{corollary}

\textbf{向量场与子流形}

若 $S\subseteq M$ 是浸入或嵌入子流形（带边或不带边），则 $M$ 上的向量场 $X$ 不一定限制为 $S$ 上的向量场，因为在点 $p\in S$ 处 $X_p$ 可能不在子空间 $T_pS\subseteq T_pM$ 中。给定点 $p\in S$，若 $X_p\in T_pS\subseteq T_pM$，则称 $M$ 上的向量场 $X$ 在 $p$ 处\textbf{与 $S$ 相切}。若它在 $S$ 的每一点处都与 $S$ 相切，则称它\textbf{与 $S$ 相切}。

\begin{figure}[ht]
\centering
\includegraphics[width=0.8\textwidth]{fig-08-4.png}
\caption{切于子流形的向量场}\label{figure:ch08-vector-field-tangent}
\end{figure}

\begin{proposition}\label{proposition:ch08-22}
设 $M$ 是光滑流形，$S\subseteq M$ 是带边或不带边的嵌入子流形，$X$ 是 $M$ 上的光滑向量场。则 $X$ 与 $S$ 相切，当且仅当对每个满足 $f|_S\equiv0$ 的 $f\in C^\infty(M)$，有 $(Xf)|_S=0$。
\end{proposition}

\begin{proof}
这是\cref{proposition:ch05-37}的直接推论。
\end{proof}

假设 $S\subseteq M$ 是带边或不带边的浸入子流形，$Y$ 是 $M$ 上的光滑向量场。若存在 $X\in\mathfrak{X}(S)$ 与 $Y$ 是 $\iota$-相关的（其中 $\iota\colon S\hookrightarrow M$ 是包含映射），则显然 $Y$ 与 $S$ 相切，因为对每个 $p\in S$，$Y_p=\mathrm{d}\iota_p(X_p)$ 在 $\mathrm{d}\iota_p$ 的像中。下一个命题表明逆命题成立。

\begin{proposition}\label{proposition:ch08-23}[将向量场限制到子流形]
设 $M$ 是光滑流形，$S\subseteq M$ 是带边或不带边的浸入子流形，$\iota\colon S\hookrightarrow M$ 表示包含映射。若 $Y\in\mathfrak{X}(M)$ 与 $S$ 相切，则存在 $S$ 上唯一的光滑向量场，记作 $Y|_S$，与 $Y$ 是 $\iota$-相关的。
\end{proposition}

\begin{proof}
$Y$ 与 $S$ 相切意味着由定义对每个 $p$，$Y_p$ 在 $\mathrm{d}\iota_p$ 的像中。因此，对每个 $p$ 存在向量 $X_p\in T_pS$ 使得 $Y_p=\mathrm{d}\iota_p(X_p)$。由于 $\mathrm{d}\iota_p$ 是单射，$X_p$ 唯一，故这将 $X$ 定义为 $S$ 上的粗向量场。若能证明 $X$ 光滑，则它就是与 $Y$ 是 $\iota$-相关的唯一向量场。只需证明它在每一点的邻域内光滑即可。

设 $p$ 是 $S$ 中任意点。由于浸入子流形（带边或不带边）局部嵌入，存在 $p$ 在 $S$ 中的邻域 $V$，它嵌入在 $M$ 中。设 $\bigl(U,(x^i)\bigr)$ 是 $M$ 中以 $p$ 为中心的 $V$ 的切片坐标卡（或边界切片坐标卡），使得 $V\cap U$ 是 $x^{k+1}=\cdots=x^n=0$（且若 $p\in\partial S$ 则 $x^k\geqslant0$）的子集，且 $(x^1,\dots,x^k)$ 构成 $S$ 在 $V\cap U$ 中的局部坐标。若在这些坐标中 $Y=Y^1\partial/\partial x^1+\cdots+Y^n\partial/\partial x^n$，则由我们的构造可知 $X$ 的坐标表示为 $Y^1\partial/\partial x^1+\cdots+Y^k\partial/\partial x^k$，它在 $V\cap U$ 上显然光滑。
\end{proof}

\end{document}